
\subsection{Traders who respond to being watched}\label{sec:hiding}

The results so far take the groups' rules as given. A group that expects a targeted action has a reason to change its rule. Sufficiency bounds what a supervisor who relies on prices can then achieve.

\begin{proposition}[Hiding caps targeting]\label{prop:hiding}
A supervisor observes prices, under any probes and for a record of any length, and then targets a given group or not. Being targeted costs the group $P>0$. The group chooses among available rules. Rule $\rho$ has a private cost $C(\rho)\ge0$ and leads to targeting with probability $\pi(\rho)$. Suppose one available rule $\rho_0$, with cost $C_0$, gives the population the kernel of a population that the supervisor targets with probability $\alpha$.
\begin{enumerate}[label=(\roman*),itemsep=1pt]
\item $\pi(\rho_0)=\alpha$, whatever the supervisor's probes, record length and decision rule.
\item Every rule $\rho^*$ that minimizes the group's expected cost $C(\rho)+P\pi(\rho)$ satisfies
\[
\pi(\rho^*)\ \le\ \alpha+\frac{C_0-C(\rho^*)}{P}\ \le\ \alpha+\frac{C_0}{P}.
\]
\item Suppose the group has two rules: its original rule, with cost zero and targeting probability $\pi_T$ on a record of length $T$, and $\rho_0$. It adopts $\rho_0$ whenever $\pi_T>\alpha+C_0/P$. If $\pi_T\to1$ and $P(1-\alpha)>C_0$, this holds for all large $T$, and the targeting probability is then $\alpha$.
\end{enumerate}
\end{proposition}

\begin{proof}
(i) By Theorem~\ref{thm:eqinv}, populations with the same kernel give the same law of the price record for every probe sequence and every insider and dealer rule. The supervisor's decision is a function of that record and of its own randomization. (ii) $C(\rho^*)+P\pi(\rho^*)\le C_0+P\alpha$ by optimality, and $C(\rho^*)\ge0$. (iii) The original rule costs $P\pi_T$ and $\rho_0$ costs $C_0+P\alpha$.
\end{proof}

\paragraph{In words.} The simplest way to hide is to adopt another group's memory. By Theorem~\ref{thm:merged} the two groups then form one atom of the merged measure, and $\alpha$ is the probability with which the supervisor targets a group that does not exist separately. The bound does not involve the amount of data. The supervisor's excess targeting probability times the penalty is at most the cost of hiding. More data make the original rule easier to detect and therefore make hiding more attractive: in part~(iii) the targeting probability falls to $\alpha$ once the record is long enough. The proposition is a revealed-preference inequality. What the paper adds is that sufficiency makes the mimicry exact, for every record and every probe. Three limits apply. The cost $C_0$ is not modelled. The proposition concerns one group and not an equilibrium among several. And a group that hides changes the kernel, and so prices; we do not say whether that helps or harms the market. A class tag is different from an inferred memory: it is attached to the account, and changing the trading rule does not remove it. With several groups, the bound holds group by group in any equilibrium: if, given the other groups' rules, group $i$ can adopt at cost $C_i$ a rule that gives the population the kernel of a population targeted with probability $\alpha_i$, then part~(ii) applies to $i$'s best response, so $i$ is targeted with probability at most $\alpha_i+C_i/P_i$. What determines the $\alpha_i$ in equilibrium is not modelled.

